\section{Related Work}
\label{sec:related}

Computational design space exploration has been studied most extensively in procedural
content generation~\cite{shaker2016procedural}, where sparse regions in level design
spaces are identified by generative search and presented to designers as novel starting
points. Liapis et al.~\cite{liapis2013sentient} demonstrate that PCA-identified sparse
regions in a space of procedurally generated game levels correspond to qualitatively
distinct and potentially interesting design directions. Our work applies the same
principle to the human-designed game corpus: we identify sparse regions not by
generative exploration but by analysis of what has actually been published.

Market gap analysis in board game publishing has been discussed by practitioners
(e.g., Caylus designer William Attia's public design notes) but not formally studied.
The closest academic precedent is Tschang's~\cite{tschang2005videogames} study of
novelty in video game design, which finds that successful innovations cluster at the
boundaries of established genres rather than in unexplored territory --- a hypothesis
our expert validation study is positioned to test for physical games.
